% 第五部：評價決定程序的判準（第參部分，1369–1376）。出處：細目次第三節參；英文底稿 Quellen/waldron_source.tex。
\chapter{評價決定程序的判準：論證的形式}

\qaentry{第五部之一}{為何需要一個決定程序}

\begin{qaquestion}
在四項假設之下，社會為何必須另設一個決定程序？「立法有效，除非侵害權利」為何不能充當決定規則？
\end{qaquestion}

權利分歧的雙方都需要一個共同行動的基礎，權利問題也在立法場域出現\pg{1369--1370}。待答的問題是：立法機關就權利問題所為的決議應否具決定性，抑或由司法機關事後複核\pg{1371}。作者對決定程序提出兩項形式要求\pg{1371--1372}：

\begin{enumerate}
\item 「立法機關的決定有效，除非侵害權利」不能充當決定規則，因為社會成員爭執的正是某一法律案是否侵害權利。作者引 Hobbes 為據：決定程序的運作必須平息爭議，不得重燃當初使人需要決定程序的那些爭議\pg{1371}。
\item 社會成員雖對權利有分歧，仍須共享一套關於決定程序之正當性的理論\pg{1371}；程序的設計必須獨立於它要解決的那項分歧\pg{1373}。
\end{enumerate}

對「程序本身亦有分歧」的反對意見（Christiano 的無窮後退），作者答以：不論在決定程序上是否也有分歧，都必須處理正當性問題\pg{1371--1372}。任何程序都會出錯，Wollheim 弔詭（同一公民可以同時主張某政策不應施行、又因多數選擇而應施行）是政治理論的一般問題，違憲審查同樣會作出錯誤決定\pg{1372}。

\qaentry{第五部之二}{結果面理由與程序面理由}

\begin{qaquestion}
評價決定程序的理由分為哪兩類？各自的定義與限制為何？
\end{qaquestion}

\subsection{程序面的理由}

程序面的理由問的是何人有權參與決定，與決定的結果無涉。作者以父母與路人之例說明：孩子犯錯該不該管教，由父母決定，街上的路人或公車上的乘客無權決定，即使我們認為路人較可能作出更好的決定\pg{1372--1373}。政治平等與投票權屬於這一類。

\subsection{結果面的理由}

作者的定義是：

\begin{quote}
Outcome-related reasons, by contrast, are reasons for designing the decision-procedure in a way that will ensure the appropriate outcome (i.e., a good, just, or right decision).\pg{1373}
\end{quote}

結果面的理由不得以特定的爭議結果為目標，否則就是乞題。作者接受一種節制的版本：選擇最可能觸及權利真相的程序，不論真相為何\pg{1373}；他在此引 Kavanagh\pg{1373--1374}。

\subsection{兩類理由的規範性質}

Raz 以結果面理由為唯一值得考慮的理由，作者稱之為教條主義，不採\pg{1374}。作者認為兩類理由都帶有義務論上的急迫性；Dworkin 把參與的價值說成「參與上的後果」，作者認為這嚴重低估了參與權的觀念\pg{1375}。

\qaentry{第五部之三}{第參部分預告的結論}

\begin{qaquestion}
兩類理由權衡下來，作者預告的結論是什麼？
\end{qaquestion}

作者以兩種問法設定比較的問題，其一取自 Michelman（註 76），其二為作者將其前後倒置而成，接著寫：

\begin{quote}
What I will argue, in Part IV, is that the outcome-related reasons are at best inconclusive. They are important, but they do not (as is commonly thought) establish anything like a clear case for judicial review. The process-related reasons, however, are quite one-sided. They operate mainly to discredit judicial review while leaving legislative decisionmaking unscathed.\pg{1375}
\end{quote}

\begin{qafigure}{兩類理由對兩個機關的檢驗結果（1375--1393）。\label{fig:reasons}}
% 兩類理由 × 兩個機關（1375、1377–1393）。三格灰階，決定勝負的一格用強調色。
\begin{tikzpicture}[qa,cell/.style={nbox,text width=52mm,minimum height=17mm}]
  \node[hdr] at (3.15,1.25) {結果面的理由（第肆部分）};
  \node[hdr] at (8.85,1.25) {程序面的理由（第伍部分）};
  \node[hdr,anchor=east] at (0.15,0) {立法機關};
  \node[hdr,anchor=east] at (0.15,-2.0) {法院};
  \node[cell] at (3.15,0) {代表結構使受影響者的處境得以體察（1377--1378）；議會辯論可直接處理道德問題（1384--1385）};
  \node[cell] at (8.85,0) {公平選舉與多數決，給予平等的發言權與平等的決定權威（1387--1389）};
  \node[cell] at (3.15,-2.0) {個案取向、文本取向、說明理由三項優點，經檢驗不成立或兩方均有（1379--1386）};
  \node[cell,kbox] at (8.85,-2.0) {法官的民主資格遠不及議員；大法官之間的簡單多數決欠缺公平性的論證（1391--1392）};
  \node[lbl,anchor=north] at (3.15,-3.05) {至多未有定論};
  \node[lbl,anchor=north] at (8.85,-3.05) {一面倒，不利於違憲審查};
\end{tikzpicture}

\end{qafigure}

\begin{readernote}
這張表是全文的骨架：第肆部分填左欄，第伍部分填右欄。導讀時先放這張表，再分兩部講，聽者比較容易跟上。
\end{readernote}
